\documentclass[11pt]{article}

\usepackage[margin=1.15in]{geometry}
\usepackage{amsmath,amssymb,amsthm,mathtools}
\usepackage{hyperref}
\usepackage{enumitem}
\usepackage{xcolor}
\usepackage{microtype}
\usepackage{booktabs}

\hypersetup{
  colorlinks=true,
  linkcolor=blue!45!black,
  citecolor=blue!45!black,
  urlcolor=blue!45!black
}

\newtheorem{theorem}{Theorem}
\newtheorem{lemma}[theorem]{Lemma}
\newtheorem{proposition}[theorem]{Proposition}
\newtheorem{corollary}[theorem]{Corollary}
\newtheorem{definition}[theorem]{Definition}
\newtheorem{remark}[theorem]{Remark}

\newcommand{\R}{\mathbb{R}}
\newcommand{\Q}{\mathbb{Q}}
\newcommand{\Z}{\mathbb{Z}}
\newcommand{\F}{\mathcal F}
\newcommand{\Pcal}{\mathcal P}
\newcommand{\dotp}{\mathbin{\cdot}}
\newcommand{\id}{\operatorname{id}}
\newcommand{\relint}{\operatorname{relint}}

\title{The Cuboctahedron Omnihedral Billiard Problem:\\
A Lean-Checked Holonomy Normal-Form Reduction}
\author{Zach Conn}
\date{\today}

\begin{document}

\maketitle

\begin{abstract}
We describe the current Lean 4 formalization strategy for proving that the
cuboctahedron has no nonsingular periodic billiard orbit visiting each face
exactly once.  The geometric part is classical: unfold the billiard path into
a straight line through reflected copies of the polyhedron.  The new proof
engineering pivot is a holonomy normal form for opposite-face pair words.  It
tracks the parity of the three square-pair reflections and a projective
triangular shadow.  In Lean, for every valid started pair word, the linear
return map is the identity if and only if the reduced triangular shadow is
empty.  The nonempty-shadow branch is the non-translation branch; the
empty-shadow branch is the translation branch.

The final theorem is reduced to semantic finite coverage over this split.
External programs may search for compression patterns and propose generated
evidence, but the trusted proof boundary remains Lean: exact rational or
integer arithmetic, no floating point, no numerical tolerances, and no
unchecked certificate claims.  At the time of this draft the public Lean
theorem is conditional on complete generated coverage; the latest strategy is
to complete that coverage by semantic holonomy-state families rather than raw
rank or rank-mask tables.
\end{abstract}

\paragraph{AI usage.}
The accompanying Lean development and this draft were produced with substantial
assistance from OpenAI Codex.  The proof artifact is the Lean repository; the
AI transcript and generated search diagnostics are not trusted proof.

\section{Introduction}

A billiard trajectory in a polyhedron moves along straight line segments and
reflects specularly at faces.  We exclude singular impacts at edges and
vertices.  The question studied here is whether the cuboctahedron admits a
periodic trajectory that hits all fourteen faces exactly once in one period.

\begin{definition}
Let \(P\subset\R^3\) be a polyhedron with face set \(\F\).  A periodic
billiard orbit in \(P\) is \emph{omnihedral} if, during one period, it meets
every face of \(P\) exactly once.  It is \emph{nonsingular} if every impact
point lies in the relative interior of the face it meets.
\end{definition}

\begin{theorem}[Target theorem]
The cuboctahedron has no nonsingular periodic billiard orbit which meets each
of its fourteen faces exactly once in one period.
\end{theorem}

This is a finite but very large exact obstruction problem.  The current
formalization proves the geometric bridge from complete finite coverage to the
target theorem, and it now contains the holonomy normal-form theorem that
separates identity-linear words from nonidentity-linear words without trusting
an external scan.  The remaining work is to emit and check complete semantic
coverage for the two branches.

The broader billiard context is familiar.  In rational polygonal billiards the
unfolding construction leads to translation surfaces and a rich theory of
periodic directions; see for instance Masur \cite{masur} and Boshernitzan,
Galperin, Kruger, and Troubetzkoy \cite{boshernitzan}.  In polyhedral
billiards the singular set and unfolded geometry are less rigidly organized.
Work such as Cekic, Georgiev, and Mukherjee \cite{cekic-georgiev-mukherjee}
illustrates that even controlling trajectories away from singularities can be
subtle.  The omnihedral condition imposes a different constraint: a minimal
face tour.  Bedaride and Rao prove that regular simplices admit such orbits
\cite{bedaride-rao}; the cuboctahedron is a natural highly symmetric test case
with two face types.

\paragraph{Proof artifact.}
The human-readable argument is intentionally secondary to the Lean
formalization at
\[
  \url{https://github.com/zpconn/cuboctahedron}.
\]
The paper explains the mathematical and proof-engineering architecture.  The
authoritative finite verification must be a Lean theorem whose statement says
that every rank and every relevant sign mask is killed by a checked
obstruction.

\section{The Exact Cuboctahedron Model}

We use the rational coordinate model
\[
P=\{(x,y,z)\in\R^3:
 |x|\le 1,\ |y|\le 1,\ |z|\le 1,\
 \pm x\pm y\pm z\le 2\}.
\]
The square faces are
\[
X_\pm:\ \pm x=1,\qquad
Y_\pm:\ \pm y=1,\qquad
Z_\pm:\ \pm z=1,
\]
and the triangular faces are
\[
T_{\epsilon_1\epsilon_2\epsilon_3}:
\epsilon_1x+\epsilon_2y+\epsilon_3z=2,
\qquad \epsilon_i\in\{\pm1\}.
\]
All supporting planes have rational normal and rational offset.  Reflection
in the plane \(n_F\dotp p=c_F\) is
\[
r_F(p)=p-2\,\frac{n_F\dotp p-c_F}{n_F\dotp n_F}\,n_F,
\]
so every face reflection has rational affine coefficients.

The fourteen faces form seven opposite-face pairs:
\[
x,\ y,\ z,\ d_{111},\ d_{11\bar 1},\ d_{1\bar 1 1},\
d_{\bar 1 11}.
\]
Opposite faces have the same linear reflection part and differ only in the
affine translation term.  This is why the proof first studies opposite-pair
words and then, in the translation case, adds the remaining sign choices.

\section{Unfolding and the Finite Language}

Let \(F_0,\ldots,F_{13}\) be a proposed ordered list of faces.  Reflecting the
polyhedron successively across those faces turns a broken billiard path into a
straight line in a reflected development.  The product of the fourteen affine
reflections is an affine return map
\[
A(p)=Mp+b,\qquad M\in O(3).
\]
In the cuboctahedron model above, \(M\) and \(b\) are rational.

\begin{lemma}[Unfolding criterion]
If a nonsingular billiard orbit meets \(F_0,\ldots,F_{13}\) in that cyclic
order, then there is an unfolded line \(\ell(t)=p+tv\) crossing the developed
faces in the corresponding order, with all crossings in relative interiors,
and the return map carries the line to itself with \(Mv=v\).  Conversely, such
an unfolded witness folds to a nonsingular billiard orbit with the prescribed
itinerary.
\end{lemma}

\begin{proof}
At each impact, reflect the next copy of the polyhedron rather than the
velocity vector.  The reflected outgoing segment becomes the continuation of
the incoming segment.  Iterating turns the billiard path into one line.
Nonsingularity is exactly the condition that the crossings remain in relative
face interiors.  Periodicity identifies the developed line with its image
under the affine return map, giving \(Mv=v\).  The construction is reversible
because reflecting the copies back restores the equal-angle law.
\end{proof}

An omnihedral period hits \(X_+\) once, so cyclic reindexing lets us assume the
first face is \(X_+\).  After fixing this start, the remaining thirteen
opposite-face entries have multiplicities
\[
x:1,\qquad y:2,\qquad z:2,\qquad
d_{111},d_{11\bar 1},d_{1\bar 1 1},d_{\bar 1 11}:2.
\]
The exact Lean enumeration proves
\[
  \#\{\hbox{valid started pair words}\}
  =\frac{13!}{2^6}=97{,}297{,}200.
\]
Ranks in \(\{0,\ldots,97{,}297{,}199\}\) are used as addresses for this
bijection, not as the main mathematical compression coordinate.

\begin{definition}
A \emph{pair word} is a length-13 word in the seven opposite-face pairs with
the multiplicities above.  A \emph{sign mask} is the remaining choice of one
face from each non-start opposite pair.  In the identity-linear branch there
are six independent sign bits, hence \(64\) masks.
\end{definition}

For a pair word \(w\), let \(M(w)\) be the linear part of the return map for
the all-positive lift.  It depends only on the opposite-face pairs.  The
product-order convention in the Lean development is
\[
  R(w_0)R(w_1)\cdots R(w_{12})R(x),
\]
where the final \(R(x)\) is the started \(X_+\) pair reflection.

\section{The Holonomy Normal Form}

The latest strategy replaces the old rank-interval-first split by a small
exact normal form for the linear holonomy.  The rank still gives exhaustive
enumeration, but the proof tries to compress by geometric state.

\subsection{Square parity and triangular shadow}

The square-pair reflections \(x,y,z\) are diagonal sign changes.  Their
accumulated contribution is a parity state
\[
  p=(p_x,p_y,p_z)\in(\Z/2\Z)^3.
\]
The triangular pair reflections are represented projectively by four letters
\[
  111,\qquad 11\bar 1,\qquad 1\bar 1 1,\qquad \bar 1 11,
\]
where changing the sign of all three coordinates gives the same projective
normal.

When a triangular reflection is scanned after square parity \(p\), it is
conjugated by that parity and appended to a triangular shadow.  In the Lean
definitions this action toggles the two projective signs by
\[
  y\hbox{-sign}\mapsto y\hbox{-sign}+p_x+p_y,\qquad
  z\hbox{-sign}\mapsto z\hbox{-sign}+p_x+p_z
\]
modulo \(2\).  Adjacent equal triangular letters cancel, because the
corresponding reflection squared is the identity.  The result is a reduced
stack, denoted here by \(\operatorname{red}(w)\).

\begin{definition}
The \emph{reduced triangular shadow} of a started pair word \(w\) is obtained
by scanning
\[
  w_0,w_1,\ldots,w_{12},x,
\]
updating square parity on square-pair letters, appending the parity-conjugated
triangular letter on triangular-pair letters, and cancelling adjacent equal
triangular letters.
\end{definition}

\subsection{The Lean classifier theorem}

The central new checkpoint is the following theorem, formalized in the
normal-form modules.

\begin{theorem}[Identity classifier]
For every valid started pair word \(w\),
\[
  M(w)=I
  \quad\Longleftrightarrow\quad
  \operatorname{red}(w)=[].
\]
\end{theorem}

\begin{proof}[Proof idea]
The proof has four exact pieces.

First, valid pair-word counts imply that the final square parity is trivial:
the scan sees \(x,y,z\) each an even number of times after including the final
started \(x\) factor.

Second, the scanner satisfies a product invariant.  If \(S_p\) is the diagonal
matrix for square parity \(p\), and \(H_t\) is the triangular reflection for a
letter \(t\), then
\[
  S_pH_t=H_{p\cdot t}S_p.
\]
Thus the original linear product is the triangular shadow product followed by
the final square parity matrix.  Since the final square parity is \(I\), the
linear product equals the triangular shadow product.

Third, cancelling adjacent equal triangular letters preserves the triangular
product.  Therefore the triangular shadow and the reduced triangular shadow
have the same product.

Fourth, a nonempty reduced triangular word cannot have product \(I\).  For a
triangular normal \(n\in\{\pm1\}^3\), write
\[
  H_n=I-\frac{2}{3}nn^T,\qquad
  K_n=3H_n=3I-2nn^T\in M_3(\Z).
\]
Modulo \(3\), one has \(K_n\equiv nn^T\).  A product of rank-one matrices
\[
  n_1n_1^T n_2n_2^T\cdots n_mn_m^T
\]
is a nonzero scalar multiple of \(n_1n_m^T\) over \(\mathbb F_3\), provided
adjacent projective letters are distinct; the relevant adjacent dot products
are nonzero in \(\mathbb F_3\).  Hence a nonempty reduced triangular product
is nonzero modulo \(3\) after integer scaling.  If the rational product were
the identity, the scaled integer product would be \(3^m I\), which is zero
modulo \(3\) for \(m>0\), a contradiction.
\end{proof}

The external full scan that motivated this formalization found
\[
\begin{array}{lr}
\text{empty reduced shadow} & 2{,}468{,}088,\\
\text{nonempty reduced shadow} & 94{,}829{,}112.
\end{array}
\]
These counts were not trusted proof.  Their role was to validate the pivot
before formalizing the classifier.  The trusted branch split is now the Lean
theorem above.

\section{The Two Obstruction Branches}

\subsection{Nonempty shadow: non-translation holonomy}

If \(\operatorname{red}(w)\ne[]\), then \(M(w)\ne I\).  A feasible unfolded
line must have direction fixed by \(M(w)\), and the affine return map must
carry the line to itself.  This forces a linear-algebraic axis condition.

\begin{lemma}[Forced-axis condition]
Let \(A(p)=Mp+b\) be the affine return map for a started itinerary with
\(M\ne I\).  If the itinerary has an unfolded periodic line witness, then the
line direction is a nonzero fixed vector of \(M\), and the starting point lies
on the affine solution set determined by \(A\) and that fixed direction.
\end{lemma}

\begin{proof}
The unfolding criterion gives \(Mv=v\).  The condition that \(A\) maps the
line \(p+\R v\) to itself is equivalent to
\[
  (I-M)p=b-\lambda v
\]
for some \(\lambda\in\R\).  This is a linear system for the starting point and
\(\lambda\).  The certificates used in the nonidentity branch are exact
rational witnesses that this forced geometry is incompatible with the required
face order or relative-interior conditions.
\end{proof}

The old approach tried to cover nonidentity ranks directly.  The active plan
uses nonempty reduced shadows, primitive-axis signatures, and forced-axis
failure signatures as semantic coordinates.  The public statement remains
small:
\[
  \text{for every rank }r,\quad
  M(\operatorname{unrank}(r))\ne I
  \Rightarrow \hbox{no unfolded feasible witness.}
\]
In Lean this is the semantic predicate
\(\texttt{NonIdentityRankKilled}\).  Generated files should prove such killed
statements through reusable families, not by exporting one raw certificate per
rank.

\subsection{Empty shadow: translation holonomy}

If \(\operatorname{red}(w)=[]\), the linear part is the identity and the return
map is a translation
\[
  A(p)=p+b.
\]
For a nondegenerate translation case, an invariant unfolded line has direction
parallel to \(b\).  With the started point written as
\[
  p_0=(1,y,z),
\]
the face-crossing and relative-interior conditions become strict rational
linear inequalities in the two real variables \(y,z\), once a sign mask is
fixed.

For the \(i\)-th unfolded face plane \(u_i\dotp p=d_i\), the crossing time is
\[
  t_i=\frac{d_i-u_i\dotp p_0}{u_i\dotp b}.
\]
The denominator \(u_i\dotp b\) is an exact rational expression depending on
the pair word and sign mask.

\begin{lemma}[Good direction]
Every feasible translation-type unfolded witness satisfies
\[
  0<u_i\dotp b
\]
for every internal impact \(i\).
\end{lemma}

\begin{proof}
In the chosen orientation the line must cross each successive unfolded face in
forward time and from the correct side.  A nonpositive internal denominator
prevents that crossing from occurring with the required orientation.  The Lean
formalization packages this as the theorem
\(\texttt{goodDirection\_of\_translation\_feasible\_at\_rank}\).
\end{proof}

This theorem is one of the main proof-engineering simplifications.  Generated
translation evidence should not mention masks that fail GoodDirection.  Those
masks are impossible by a hand-written semantic theorem.  The generated
translation task is only:
\[
  \text{identity-linear rank }r,\quad
  \text{mask }m,\quad
  \text{GoodDirection}(r,m)
  \Longrightarrow
  \hbox{case killed.}
\]

\subsection{Strict Farkas obstructions}

For a fixed translation rank and GoodDirection mask, the remaining constraints
are strict rational inequalities
\[
  a_jy+c_jz+e_j>0,\qquad a_j,c_j,e_j\in\Q.
\]

\begin{lemma}[Strict two-variable Farkas certificate]
Suppose there are rational multipliers \(q_j\ge0\), not all zero, such that
\[
  \sum_j q_ja_j=0,\qquad
  \sum_j q_jc_j=0,\qquad
  \sum_j q_je_j\le0.
\]
Then the strict system is infeasible.
\end{lemma}

\begin{proof}
If \(y,z\) satisfied all inequalities, multiplying by \(q_j\) and summing
would give
\[
  0<
  \sum_jq_j(a_jy+c_jz+e_j)
  =
  \Big(\sum_jq_ja_j\Big)y+
  \Big(\sum_jq_jc_j\Big)z+
  \sum_jq_je_j
  \le0,
\]
a contradiction.
\end{proof}

The current translation pipeline favors \emph{source-indexed} Farkas proofs.
A row is not identified merely by a large generated number; it is identified
by its geometric source, such as a start-face interior condition, an ordering
condition, or an impact-face interior condition.  Many survivors are killed by
two-source contradictions.  The proof object records the two sources and the
exact rational relation between the corresponding rows.

\subsection{Walsh and integer-normalized evidence}

The six sign bits in the translation branch may be represented by
\[
  s_1,\ldots,s_6\in\{\pm1\}.
\]
For a fixed pair word, internal denominators are degree-at-most-two Walsh
polynomials in these signs.  This follows because both the unfolded normal and
the translation vector are affine sign expressions, and their dot product is
quadratic.

Walsh subcube bounds remain useful as local proof evidence and diagnostics:
if a denominator, or a nonnegative weighted sum of denominators, is
nonpositive on a sign subcube, then every mask in that subcube violates
GoodDirection.  The active completion path, however, no longer tries to finish
the proof by a raw mask-subcube tiling.  Several exact bounded audits showed
that bad-direction masks are too irregular in rank-mask coordinates.  The
preferred generated backend now uses source-index/state classifiers,
row-producer families, and integer or homogeneous data where possible, with
small rational interpretation bridges proved once in Lean.

\section{Generated Coverage as a Semantic Theorem}

The trusted finite theorem should not expose giant arrays of certificates.
The public generated surface is semantic:
\[
\begin{array}{ll}
\texttt{NonIdentityRankKilled}(r) &
  \text{rank }r\text{ is impossible when }M(r)\ne I,\\
\texttt{TranslationCaseKilled}(r,m) &
  \text{mask }m\text{ at rank }r\text{ is impossible when }M(r)=I.
\end{array}
\]
The structure \(\texttt{SemanticExhaustiveGeneratedCoverage}\) packages:
\begin{enumerate}[leftmargin=2em]
\item every pair rank is covered by the enumeration address space;
\item every sign mask is in the finite sign-mask address space;
\item every nonidentity rank is killed;
\item every translation rank and mask is killed.
\end{enumerate}

For the translation branch, the newest public route is all-Good coverage.  It
asks generated evidence to prove only the GoodDirection survivors and then
uses the theorem ``translation feasible implies GoodDirection'' to recover
ordinary translation coverage.  In repository terms, the route is
\[
\texttt{PublicAllGoodSemanticCoverageIntervals}
\longrightarrow
\texttt{semanticGeneratedCoverageOfAllGoodIntervals}
\longrightarrow
\texttt{SemanticExhaustiveGeneratedCoverage}.
\]
Although the public assembly may still be phrased over intervals
\([lo,hi)\), those intervals are an assembly device.  The mathematical
compression coordinate is now holonomy state and semantic family data, not
lexicographic rank adjacency.

\section{The Conditional Lean Bridge}

The hand-written Lean proof already connects complete generated coverage to
the real billiard theorem.  The current public bridge has the shape
\begin{verbatim}
theorem Cuboctahedron.conditional_cuboctahedron_no_omnihedral
    (coverage : ExhaustiveGeneratedCoverage) :
    Not (exists o : BilliardOrbit14,
      o.Nonsingular /\ o.Periodic /\ o.TouchesEachFaceExactlyOnce)
\end{verbatim}
There is also a semantic coverage route that first converts semantic killed
statements into the older coverage structure expected by this bridge.

\begin{proposition}[Conditional completion]
If Lean is given exhaustive semantic generated coverage for all started pair
words and all relevant sign masks, then the target theorem follows.
\end{proposition}

\begin{proof}
Assume a nonsingular omnihedral orbit exists.  Cyclically reindex it to start
on \(X_+\).  The unfolding theorem gives an unfolded feasible started
sequence.  The itinerary theorem converts that sequence to a valid pair word,
hence to a rank.  The holonomy classifier splits the rank by whether the
reduced shadow is empty.

If the reduced shadow is nonempty, then the nonidentity semantic coverage
kills the rank.  If it is empty, the sequence determines one of the \(64\)
sign masks.  Feasibility implies GoodDirection; all-Good semantic translation
coverage kills that survivor, and the GoodDirection adapter gives ordinary
translation impossibility.  Both branches contradict the unfolded feasible
witness.
\end{proof}

This proposition is the logical end of the proof.  The remaining task is not a
new geometric idea, but a complete Lean-checked emission of the semantic
coverage theorem at a scale that compiles.

\section{Current Status and Validation}

The current formalization has the following checked structural components.
\begin{enumerate}[leftmargin=2em]
\item Exact cuboctahedron faces, reflections, unfolded feasibility, and the
  billiard-to-unfolding bridge.
\item The rank/unrank enumeration of all \(97{,}297{,}200\) started pair
  words.
\item Certificate soundness for nonidentity forced-axis obstructions,
  translation linear-inequality obstructions, and strict two-variable Farkas
  certificates.
\item The translation GoodDirection theorem, so bad-direction masks do not
  become generated proof obligations.
\item Semantic generated-coverage predicates and adapters, including the
  all-Good translation route.
\item The holonomy normal-form classifier:
  \[
    M(w)=I\quad\Longleftrightarrow\quad \operatorname{red}(w)=[].
  \]
\end{enumerate}

The unconditional final theorem is not yet the public Lean theorem.  The
public evidence marker is bounded rather than exhaustive.  The next proof
engineering step is to generate and check complete semantic coverage over the
full rank range, using the normal-form split and state-level families.

The expected validation commands for the checked Lean artifact are:
\begin{verbatim}
lake build
grep -R "sorry\|admit\|axiom\|native_decide\|unsafe" Cuboctahedron || true
lean Cuboctahedron/PrintAxioms.lean
\end{verbatim}
Until the unconditional theorem is installed, the conditional axiom audit is
the relevant public check:
\begin{verbatim}
lake env lean Cuboctahedron/PrintConditionalAxioms.lean
\end{verbatim}

\section{Why the Strategy Changed}

The proof began as a direct finite enumeration: produce one obstruction for
each rank or rank-mask case and ask Lean to check all of them.  That is
mathematically sound in principle and unusable in practice.  The source files,
kernel reductions, rational arithmetic, and memory usage grew too quickly.
Packed blobs, broad Boolean checkers, raw singleton rank repairs, and blind
global rank sweeps all failed some combination of source-size, wall-time, or
RAM gates.

The holonomy pivot changes the shape of the problem.  It proves the identity
split once, in hand-written Lean, and then uses the reduced shadow as a
geometric state rather than rediscovering \(M=I\) by generated computation.
The translation side similarly uses GoodDirection as a theorem rather than
generating evidence for masks that already violate a necessary condition.

Thus the current proof-search principle is:
\[
\text{find semantic states that kill many cases, then make Lean check the
state theorem.}
\]
Ranks and intervals remain essential for exhaustiveness, but they are now the
addressing layer, not the main source of compression.

\section{Conclusion}

The cuboctahedron problem has a compact geometric reduction and a large exact
finite core.  The latest Lean strategy isolates the finite core behind a
holonomy normal form and semantic coverage predicates.  The main structural
advance is the Lean-checked theorem that, for every valid started pair word,
identity holonomy is exactly empty reduced triangular shadow.  This turns the
old translation/nontranslation split into a small theorem rather than a
trusted external classification.

Completing the proof now means completing the generated semantic coverage
over the two normal-form branches.  Once that coverage is supplied and checked
by Lean, the existing conditional bridge yields the nonsingular omnihedral
nonexistence theorem for the cuboctahedron.

\begin{thebibliography}{9}

\bibitem{tabachnikov}
S. Tabachnikov,
\emph{Geometry and Billiards},
Student Mathematical Library, Vol. 30,
American Mathematical Society, 2005.

\bibitem{masur}
H. Masur,
Closed trajectories for quadratic differentials with an application to
billiards,
\emph{Duke Mathematical Journal} 53 (1986), no. 2, 307--314.

\bibitem{boshernitzan}
M. Boshernitzan, G. Galperin, T. Kruger, and S. Troubetzkoy,
Some remarks on periodic billiard orbits in rational polygons,
\emph{Commentarii Mathematici Helvetici} 73 (1998), 177--194.

\bibitem{bedaride-rao}
N. Bedaride and M. Rao,
Regular simplices and periodic billiard orbits,
\emph{Journal of Modern Dynamics} 6 (2012), no. 2, 253--265.

\bibitem{cekic-georgiev-mukherjee}
M. Cekic, B. Georgiev, and M. Mukherjee,
Polyhedral billiards, eigenfunction concentration and almost periodic control,
\emph{Communications in Mathematical Physics} 383 (2021), 377--405.

\bibitem{lean}
L. de Moura and S. Ullrich,
The Lean 4 theorem prover and programming language,
in \emph{Automated Deduction -- CADE 28},
Lecture Notes in Computer Science 12699, Springer, 2021, pp. 625--635.

\end{thebibliography}

\end{document}
